\documentclass[10pt,reqno]{amsart}

\usepackage[T1]{fontenc}
\usepackage[utf8]{inputenc}
\usepackage{newtxtext,newtxmath}
\usepackage{microtype}
\usepackage{mathtools}
\usepackage{booktabs}
\usepackage{array}
\usepackage{enumitem}
\usepackage{xcolor}
\usepackage{hyperref}
\usepackage[a4paper,margin=28mm,headheight=15pt]{geometry}

\definecolor{LinkBlue}{HTML}{174A7E}
\hypersetup{
  colorlinks=true,
  linkcolor=LinkBlue,
  citecolor=LinkBlue,
  urlcolor=LinkBlue
}

\setlength{\parskip}{2pt}
\setlist[itemize]{leftmargin=6mm,itemsep=1.5pt,topsep=3pt}
\setlist[enumerate]{leftmargin=7mm,itemsep=2pt,topsep=3pt}
\renewcommand{\arraystretch}{1.15}
\emergencystretch=2em

\newtheorem{theorem}{Theorem}[section]
\newtheorem{proposition}[theorem]{Proposition}
\newtheorem{lemma}[theorem]{Lemma}
\newtheorem{corollary}[theorem]{Corollary}
\theoremstyle{definition}
\newtheorem{definition}[theorem]{Definition}
\newtheorem{remark}[theorem]{Remark}

\title[Quadratic Reduction and Scalar-Span Obstruction for Erd\H{o}s--Ulam Rational Distances]
{Bottom-Row Quadratic Reduction and Scalar-Span Obstruction on Rank-Four Parabolic Erd\H{o}s--Ulam Configurations}
\author{Argus Open Research Project}
\date{September 5, 2026}
\subjclass[2020]{Primary 52C10; Secondary 11G05, 14G05}
\begin{document}
\begin{abstract}
The Erd\H{o}s--Ulam problem asks whether there exists a dense subset of $\mathbb{R}^2$ all of whose pairwise Euclidean distances are rational. In the Gaussian rank-four parabolic parameterization at scale $\lambda=27$, we prove an exact quadratic reduction of the bottom distance row and a scalar-span rank-deficiency obstruction ruling out the rank-deficient non-star sub-branch.
\end{abstract}
\maketitle
\section{Main Algebraic Reduction}
\begin{proposition}
On the Gaussian rank-four parabolic locus at parameter $\lambda=27$, clearing the first three base-point rational-distance squares reduces the fourth distance condition to a binary quadratic norm equation whose scalar-span specialization has rank at most~$2$ and admits only collinear/star-degenerate rational solutions.
\end{proposition}
\end{document}
